% from ega0-15-fr.tex
\oldpage[0\textsubscript{IV-1}]{18}
est exacte ($0_{\mathrm I}$, 6.1.2). Cela montre que l'homothétie de rapport
$f'_1$ dans $M'$ est injective.

Posons ensuite, pour $1\leq i\leq n-1$,
\[
  M_i=M/\left(\sum_{j=1}^i f_jM\right),\qquad
  M'_i=M'/\left(\sum_{j=1}^i f'_jM'\right)~;
\]
on a $M'_i=M_i\otimes_A A'$, $M_i/f_{i+1}M_i=M_{i+1}$,
$M'_i/f'_{i+1}M'_i=M'_{i+1}$~; il suffit de remplacer dans le raisonnement
précédent $M$ et $f_1$ par $M_i$ et $f_{i+1}$ pour conclure que l'homothétie
de rapport $f'_{i+1}$ dans $M'_i$ est injective, en raison de la platitude de
$M_i/f_{i+1}M_i$. La dernière assertion résulte de ($0_{\mathrm I}$, 6.2.1).
\end{proof}

\begin{proposition}[15.1.16]
\label{0.15.1.16-fr}
Soient $A$, $B$ deux anneaux locaux noethériens, $k$ le corps résiduel de
$A$, $\varphi:A\to B$ un homomorphisme local, $M$ un $B$-module de type fini.
Soit $(f_i)_{1\leq i\leq n}$ une suite d'éléments de l'idéal maximal de $B$,
et pour tout $i$, soit $g_i$ l'image de $f_i$ dans $B\otimes_A k$. Les
conditions suivantes sont équivalentes~:
\begin{enumerate}
  \item[\rm(a)] La suite $(f_i)$ est $M$-régulière et les quotients
  \[
    M_i=M/\left(\sum_{j=1}^i f_jM\right)
  \]
  sont des $A$-modules plats pour $1\leq i\leq n$.
  \item[\rm(b)] La suite $(f_i)$ est $M$-régulière et le $A$-module $M_n$ est
  plat.
  \item[\rm(c)] Le $A$-module $M$ est plat et la suite
  $(g_i)_{1\leq i\leq n}$ est $(M\otimes_A k)$-régulière.
  \item[\rm(d)] Le $A$-module $M$ est plat, et pour tout homomorphisme
  d'anneaux $\rho:A\to A'$, si on pose $B'=B\otimes_A A'$,
  $M'=M\otimes_A A'$, $f'_i=f_i\otimes1$ ($1\leq i\leq n$), la suite
  $(f'_i)$ est $M'$-régulière.
\end{enumerate}
\end{proposition}

\begin{proof}
Il est clair que a) entraîne b)~; pour montrer que b) entraîne d), il suffit,
en vertu de (15.1.15), de montrer que b) entraîne que les $M_i$ sont des
$A$-modules plats pour $0\leq i\leq n-1$ (en posant $M_0=M$). Comme
$M_{i+1}=M_i/f_{i+1}M_i$, on est ramené, par récurrence descendante, à prouver
que si $h$ est un élément de l'idéal maximal de $B$ tel que l'homothétie de
rapport $h$ dans $M$ soit injective et que $M/hM$ soit un $A$-module plat,
alors $M$ est un $A$-module plat~; mais cela n'est autre que
($0_{\mathrm{III}}$, 10.2.7). Il est clair que c) est le cas particulier de
d) obtenu en prenant $A'=k$. Reste donc à prouver que c) entraîne a). Comme
$M$ est un $A$-module plat et un $B$-module de type fini, il résulte de c) et
de ($0_{\mathrm{III}}$, 10.2.4) que l'homothétie de rapport $f_1$ dans $M$
est injective et que son conoyau $M/f_1M=M_1$ est un $A$-module plat.
Raisonnons par récurrence sur $i$ et supposons démontré que $M_i$ est un
$A$-module plat~; comme l'homothétie de rapport $g_{i+1}$ dans
$M_i\otimes_A k$ est injective par hypothèse, il résulte encore de
($0_{\mathrm{III}}$, 10.2.4) que l'homothétie de rapport $f_{i+1}$ dans
$M_i$ est injective et que $M_i/f_{i+1}M_i=M_{i+1}$ est un $A$-module plat,
ce qui achève la démonstration. On notera que la démonstration du fait que c)
entraîne a) n'utilise pas le fait que les $f_i$ appartiennent à l'idéal
maximal de $B$.
\end{proof}

\begin{corollary}[15.1.17]
\label{0.15.1.17-fr}
Les hypothèses sur $A$ et $B$ étant celles de (15.1.16), soient $M$ un
$B$-module de type fini, plat sur $A$, $R$ un sous-$B$-module de $M$ tel que
$N=M/R$ soit un $A$-module plat. Soit $(f_i)_{1\leq i\leq n}$ une suite
d'éléments de l'idéal maximal de $B$ ayant les propriétés suivantes~:
\begin{enumerate}
  \item[(i)] $f_iM\subset R$ pour $1\leq i\leq n$.
  \item[(ii)] Si $g_i$ est l'image de $f_i$ dans $B\otimes_A k$, la suite
  $(g_i)_{1\leq i\leq n}$ est $(M\otimes_A k)$-régulière.
  \item[(iii)] La somme des sous-$(B\otimes_A k)$-modules
  $g_i(M\otimes_A k)$ de $M\otimes_A k$ est l'image canonique de
  $R\otimes_A k$ dans $M\otimes_A k$ (égale au noyau de l'homomorphisme
  canonique $M\otimes_A k\to N\otimes_A k$).
\end{enumerate}
Alors la suite $(f_i)$ est $M$-régulière et l'on a
$R=\sum_{i=1}^n f_iM$.
\end{corollary}

\begin{proof}
